%!TEX TS-program = xelatex
%!TEX encoding = UTF-8 Unicode
% Ogolne CV do portfolio - IT / programowanie / audyt / QA. Bez dopasowania do konkretnej oferty. Zob. changes.md.
\documentclass[11pt,a4paper]{awesome-cv}
\hyphenpenalty=10000
\exhyphenpenalty=10000
\emergencystretch=2em
\geometry{left=1.4cm,top=.5cm,right=1.4cm,bottom=.5cm,footskip=.5cm}
\fontdir[fonts/]
\colorlet{awesome}{awesome-red}
\setbool{acvSectionColorHighlight}{true}
\renewcommand{\acvHeaderSocialSep}{\quad\textbar\quad}
\name{Marek}{Kleszcz}
\address{Katowice, Polska}
\email{contact@mkleszcz.pl}
\homepage{mkleszcz.pl}
\linkedin{marek-kleszcz}
\github{kleszczuch}
\photo[rectangle,edge,right]{profile.jpg}
\renewcommand{\acvSectionTopSkip}{1.6mm}
\renewcommand{\acvSectionContentTopSkip}{1.6mm}
\renewcommand{\acvHeaderAfterSocialSkip}{2mm}
\hypersetup{pdfauthor={Marek Kleszcz},pdftitle={Marek Kleszcz - CV}}
\renewenvironment{cvparagraph}{%
  \vspace{\acvSectionContentTopSkip}\vspace{-1mm}
  \paragraphstyle
}{\par\vspace{2mm}}
\renewenvironment{cvskills}{%
  \vspace{\acvSectionContentTopSkip}\vspace{-2mm}
  \setlength{\topsep}{0pt}
  \begin{center}
  \setlength\tabcolsep{1ex}
  \setlength{\extrarowheight}{1pt}
  \begin{tabular*}{\textwidth}{@{\extracolsep{\fill}} R{3.6cm} L{\textwidth - 3.6cm - 4ex}}
}{\end{tabular*}\end{center}}

\begin{document}
\setlength{\parskip}{2pt}
\makecvheader[C]

\cvsection{Profil}

\begin{cvparagraph}
Absolwent informatyki i student studiów magisterskich ze specjalizacją cyberbezpieczeństwo. Ponad dwa lata jako tester oprogramowania w Sperasoft: testy funkcjonalne, wydajnościowe, regresyjne i penetracyjne oraz analiza i dokumentacja defektów. Od ponad roku komercyjnie tworzę oprogramowanie dla stowarzyszenia (Next.js, React, Supabase, PostgreSQL); projekty własne w Pythonie i Flutterze. Na co dzień pracuję z agentami AI i lokalnymi modelami LLM.
\end{cvparagraph}

\cvsection{Umiejętności}

\begin{cvskills}
  \cvskill{QA}{Testy funkcjonalne, wydajnościowe i regresyjne, przypadki testowe, analiza defektów i logów, Jira, JQL}
  \cvskill{Programowanie}{Python, TypeScript, JavaScript, C, C++, C\#/.NET, Java, Dart, SQL}
  \cvskill{Web i backend}{React, Next.js, Node.js, REST API, Supabase, PostgreSQL, Firebase, Flutter}
  \cvskill{AI i dane}{PyTorch, OpenCV, MediaPipe, pandas, NumPy; lokalne LLM, fine-tuning, RAG, agenci AI}
  \cvskill{Infrastruktura}{Git, Docker, GitHub Actions (CI/CD), Cloudflare, Linux, Windows Server; TCP/IP, DNS, HTTP}
\end{cvskills}

\cvsection{Doświadczenie}

\begin{cvparagraph}
\noindent\begin{minipage}[t]{\dimexpr\linewidth-4.8cm\relax}\vspace{0pt}
\textbf{Sperasoft} - Tester oprogramowania. Analiza i dokumentacja defektów. Szukanie ich źródeł poprzez logiczne myślenie i analizę zachowania aplikacji.
\end{minipage}\hfill
\begin{minipage}[t]{4.5cm}\vspace{0pt}
\raggedleft\entrylocationstyle{Katowice}\\
\entrydatestyle{czerwiec 2024 - wrzesień 2026}
\end{minipage}
\begin{itemize}[leftmargin=2ex,nosep,topsep=1mm,label=\textbullet]
  \item Testy funkcjonalne, wydajnościowe i regresyjne aplikacji o skali enterprise. Weryfikacja poprawek.
  \item Testy penetracyjne; wykonywanie planów testowych.
  \item Zgłaszanie i śledzenie błędów w Jira; komunikacja techniczna między QA, deweloperami i kierownictwem projektu.
\end{itemize}
\vspace{2mm}
\noindent\begin{minipage}[t]{\dimexpr\linewidth-4.8cm\relax}\vspace{0pt}
\textbf{Gwiazdy Katowice DESW} - Prezes stowarzyszenia. Tworzenie strony internetowej i oprogramowania dla stowarzyszenia. Koordynacja działalności klubu i pracy instruktorów.
\end{minipage}\hfill
\begin{minipage}[t]{4.5cm}\vspace{0pt}
\raggedleft\entrylocationstyle{Katowice}\\
\entrydatestyle{2025 - obecnie}
\end{minipage}
\begin{itemize}[leftmargin=2ex,nosep,topsep=1mm,label=\textbullet]
  \item Strona klubu (Next.js) i aplikacja dla członków (React, Supabase); automatyzacje miesięcznych podsumowań.
  \item Nadzór nad rozliczeniami finansowymi i dokumentacją statutową; komunikacja z instytucjami i samorządem.
  \item Organizacja wydarzenia szermierczego pod patronatem TAURON Parku Śląskiego, z uczestnikami z całej Polski.
\end{itemize}
\end{cvparagraph}

\cvsection{Projekty}

\begin{cventries}
  \cventry
    {Next.js, React, TypeScript, Vite, Supabase, PostgreSQL}
    {Gwiazdy Katowice - strona i portal stowarzyszenia}
    {\href{https://www.gwiazdykatowice.pl/pl}{\fontsize{8pt}{1em}\selectfont gwiazdykatowice.pl}}
    {}
    {\begin{cvitems}
      \item {Strona PL/EN z formularzem kontaktowym i integracjami REST po stronie serwera (Discord, Supabase).}
      \item {Portal członków: logowanie Google, PostgreSQL z kontrolą dostępu RLS/RPC, Edge Functions; w trakcie rozwoju.}
    \end{cvitems}}
  \vspace{5mm}
  \cventry
    {Dart, Flutter, Firebase Auth, Cloud Firestore}
    {HemaRef - aplikacja sędziowska HEMA}
    {\href{https://github.com/kleszczuch/HemaRef}{\fontsize{8pt}{1em}\selectfont github.com/kleszczuch/HemaRef}}
    {}
    {\begin{cvitems}
      \item {Wieloplatformowa aplikacja turniejowa (praca licencjacka): drabinka pucharowa, punktacja, kartki i kary, timer, cofanie akcji.}
    \end{cvitems}}
  \vspace{5mm}
  \cventry
    {Python, OpenCV, MediaPipe Hands, PyTorch}
    {MouseNoNeed - sterowanie gestami}
    {\href{https://github.com/kleszczuch/MouseNoNeed}{\fontsize{8pt}{1em}\selectfont github.com/kleszczuch/MouseNoNeed}}
    {}
    {\begin{cvitems}
      \item {Sterowanie myszą i klawiaturą gestami dłoni śledzonymi z kamery w czasie rzeczywistym; GUI do mapowania gestów.}
    \end{cvitems}}
\end{cventries}

\vspace{2mm}
\cvsection{Wykształcenie}

\begin{cventries}
  \cventry
    {Studia magisterskie - informatyka, cyberbezpieczeństwo}
    {Uniwersytet WSB w Katowicach}
    {Katowice}
    {październik 2026 - obecnie}
    {Studia w trakcie.}
  \vspace{2mm}
  \cventry
    {Licencjat - informatyka}
    {Uniwersytet Ekonomiczny w Katowicach}
    {Katowice}
    {2023 - 2026}
    {Specjalizacja: Algorytmy i programowanie. Praca licencjacka: aplikacja sędziowska HEMA we Flutterze i Firebase.}
\end{cventries}

\cvsection{Języki}

\begin{cvparagraph}
\textbf{Polski} - ojczysty \qquad \textbf{Angielski} - C1 \qquad \textbf{Niemiecki} - B1/B2 \qquad \textbf{Włoski} - A1/A2
\end{cvparagraph}

\vspace{2mm}
{\fontsize{8pt}{9pt}\selectfont\color{graytext}
\noindent Wyrażam zgodę na przetwarzanie moich danych osobowych zawartych w CV na potrzeby bieżącego procesu rekrutacji, zgodnie z Rozporządzeniem Parlamentu Europejskiego i Rady (UE) 2016/679 (RODO).\par}
\end{document}
